\section{\texorpdfstring{\(\GL_n \F_q\)}{GLn Fq} character theory}
\label{three}





In this section, we describe how to compute the values of all the irreducible characters of $\GL_n \F_q$.
Just as with the symmetric groups, we will index the irreducible characters of \(\GL_n \F_q\) in the same way that we have indexed its conjugacy classes.
The following is a condensed review of the topic, based on Green's work \cite{green}. The notation and language we use vary from Green's original choices. For another exposition see Macdonald \cite[{Chapter~IV}]{macD}.






\subsection{Computing the irreducible \texorpdfstring{$\GL_n \F_q$}{GLn Fq} characters} \label{char_vals}



We first introduce more notation.
For each positive integer \(d\), define \(\alpha_d : d\Z \to \Z\) by \(\alpha_d (m) = [m/d]_{q^{d}}\).
Given a polynomial \(f \in \CF (q)\), we will consider the function \(\ell_f \cdot \alpha_{\deg f}\), which is obtained by scaling \(\alpha_{\deg f}\) by the integer \(\ell_f\).



We require a process called \myemph{parabolic induction}, which we describe now.
If $\nu = (\nu_1,\ldots,\nu_\ell) \vdash n$, let $\para_\nu$ denote the \myemph{parabolic subgroup} of $\GL_{n} \F_q$ consisting of block upper-triangular matrices with block sizes $\nu_1, \ldots, \nu_\ell$. Explicitly,
\begin{equation}
\label{para_def}
\para_\nu = \left \{
\begin{bmatrix}
A_{11} & A_{12} & \cdots & A_{1 \ell} \\
0 & A_{22} & \cdots & A_{2 \ell} \\
0 & 0 & \ddots & \vdots \\
0 & 0 & 0 & A_{\ell \ell}
\end{bmatrix} \in \GL_{n} \F_q :
A_{ii} \in \GL_{\nu_i} \F_q \text{ for all } 1 \leq i \leq \ell
\right \}.
\end{equation}
For each $i \in \{1,\ldots,\ell\}$, let $\pi_i^\nu : \para_{\nu} \to \GL_{\nu_i} \F_q$ denote projection
onto the $i^\text{th}$ diagonal block:
\begin{equation}
A = \begin{bmatrix}
A_{11} & A_{12} & \cdots & A_{1 \ell} \\
0 & A_{22} & \cdots & A_{2 \ell} \\
0 & 0 & \ddots & \vdots \\
0 & 0 & 0 & A_{\ell \ell}
\end{bmatrix}
\in \para_\nu
\implies
\pi_i^\nu (A) = A_{ii}.
\end{equation}
Given arbitrary characters $\chi_i$ of $\GL_{\nu_i} \F_q$ for each $i \in \{1,\ldots,\ell\}$,
we define their parabolic induction product $\bigodot_{i = 1}^\ell \chi_i$, which is a character of $\GL_{n} \F_q$, by
\begin{equation} \label{parabolic_induction} \left(\bigodot_{i = 1}^\ell \chi_i \right) (g) = \frac{1}{\# \para_\nu} \sum_{\substack{x \in \GL_{n} \F_q \\ xgx^{-1} \in \para_\nu}} \, \, \prod_{i = 1}^{\ell} \chi_i \left ( \pi^\nu_i ( xgx^{-1} ) \right ). \end{equation}

We now define the irreducible characters of $\GL_n \F_q$ in four steps. First, we define the \myemph{Primary-support} characters, $P$.
Second, we define the \myemph{paraBolic} characters, $B$, in terms of the $P$'s.
Third, we define the \myemph{Jrreducible} characters, $J$, in terms of the $B$'s.
Finally, we define the irreducible characters \(\chi^\ulambda\) of \(\GL_n \F_q\) in terms of the $J$'s.\footnote{One might refer to this as the `PBJ' method.}
The names Primary-support, paraBolic, and Jrreducible were not used by Green.

For each $b \in \Z$ and $d \in \N$, we define the \uline{\textbf{P}}rimary-support character, $P_d^b$, of $\GL_d \F_q$ as follows:
\begin{equation}
\label{P_chars}
P_d^b (g) = \begin{dcases} \left ( \prod_{i=1}^{\ell(\mu) - 1} (1-t^i) \right ) \cdot \sum_{i=1}^{\deg h} \te(\eps_{\deg h}^{\ell_h})^{q^i b} & \text{if } \ulambda^g = h \mapsto \mu \text{ is primary}, \\ 0 & \text{otherwise.} \end{dcases}
\end{equation}
For each \(d \in \Z\), each $\nu \in \Par \setminus \{\emptyset\}$ such that \(d\) divides every part of \(\nu\), and each function $\alpha : d \Z \to \Z$, we define the para\uline{\textbf{B}}olic character, $B_\nu^\alpha$, of $\GL_{|\nu|} \F_q$ by
\begin{equation}
\label{B_chars}
B_\nu^\alpha = \bigodot_{i = 1}^{\ell(\nu)} P_{\nu_i}^{\alpha(\nu_i)}. \end{equation}
For each $f \in \CF (q)$ and $\lambda \in \Par$, we define the \uline{\textbf{J}}rreducible character, $J_f^\lambda$, of $\GL_{|\lambda| \cdot \deg f} \F_q$
by
\begin{equation}
\label{jrreducible}
J_f^\lambda
=
(-1)^{|\lambda| \cdot (\deg f - 1)} \cdot \sum_{\nu \vdash |\lambda|} \frac{1}{z_\nu} \cdot \chi^\lambda_\nu \cdot B^{\ell_f \cdot \alpha_{\deg f}}_{(\deg f) \nu}.
\end{equation}
Finally, for each index $\ulambda : \CF (q) \to \Par$ satisfying $\| \ulambda \| = n$, we define the irreducible character $\chi^\ulambda$ of $\GL_n \F_q$
by
\begin{equation} \label{grand_finale} \chi^{\ulambda} = \bigodot_{f \in \supp \ulambda} J_f^{\ulambda(f)}. \end{equation}

\begin{theorem}[Green {\cite[Theorem~14]{green}}]
\label{greens_main_thm}
For all \(n \in \N\) and prime powers~\(q\), the set $\{ \chi^\ulambda : \| \ulambda \| = n\}$ is the set of distinct, irreducible, complex characters of $\GL_n \F_q$.
\end{theorem}

Green also showed that $\odot$ is commutative and associative, so it makes sense to use an arbitrary finite indexing set for the parabolic induction product. In fact, letting $J_f^\emptyset$ denote the empty function, which is the identity element with respect to $\odot$, we can define
\begin{equation}
\chi^\ulambda = \bigodot_{f \in \CF (q)} J_f^{\ulambda(f)}.
\end{equation}

Note that we have indexed the irreducible characters in the same way that we indexed the conjugacy classes.
Moreover, we also refer to an irreducible character as \myemph{primary} if its index is primary, and we use the usual \(f \mapsto \lambda\) notation for its index.
Thus, \myemph{primary characters} are those of the form \(\chi^{f \mapsto \lambda}\) for some \(f \in \CF (q)\) and \(\lambda \in \Par\).

\begin{example}
\label{runningexample}
Using the notation from Example~\ref{conj_class_example}, we record in Table~\ref{GL3F2_char_table} the character table for $\GL_3 \F_2$.
Our choice of \(\eps\) was made such that \(f_3(\eps_3) = 0\).
The rows correspond to the characters, and the columns correspond to the conjugacy classes.
In the row labels, we write, for instance, $f_1 \mapsto (2,1)$ instead of $\chi^{f_1 \mapsto (2,1)}$ for simplicity.

\begin{table}[ht]
\caption{The character table of \(\GL_3 \F_2\), where $\zeta_7 = e^{2 \pi i / 7}$}
\centering
\begin{tabular}{ c || c | c | c | c | c | c }
 & $U_1$ & $U_2$ & $U_3$ & $E$ & $\tilde{E}$ & $C_{\mathrm{o}}$ \\ \hline \hline
$f_1 \mapsto (1,1,1)$ & $8$ & $0$ & $0$ & $1$ & $1$ & $-1$  \\
$f_1 \mapsto (2,1)$ & $6$ & $2$ & $0$ & $-1$ & $-1$ & $0$ \\
$f_1 \mapsto (3)$ & $1$ & $1$ & $1$ & $1$ & $1$ & $1$ \\
$f_3 \mapsto (1)$ & $3$ & $-1$ & $1$ & $\zeta_7 + \zeta_7^2 + \zeta_7^4$ & $\zeta_7^3 + \zeta_7^5 + \zeta_7^6$ & $0$  \\
$\tilde{f}_3 \mapsto (1)$ & $3$ & $-1$ & $1$ & $\zeta_7^3 + \zeta_7^5 + \zeta_7^6$ & $\zeta_7 + \zeta_7^2 + \zeta_7^4$ & $0$  \\
$\ulambda_{\,\mathrm{o}}$ & $7$ & $-1$ & $-1$ & $0$ & $0$ & $1$  \\
\end{tabular}
\label{GL3F2_char_table}
\end{table}
\end{example}









\subsection{Degrees of the irreducible \texorpdfstring{$\GL_n \F_q$}{GLn Fq} characters}

Green also gave an explicit formula for the degrees of the irreducible characters $\chi^\ulambda$.
For any partition $\lambda$, let $b(\lambda) = \sum_{i=1}^{\ell(\lambda)} (i-1)\lambda_i$ and define
\begin{equation}
\left [ \lambda : q \right ] = q^{b(\lambda)} \cdot \frac{\prod_{1 \leq i < j \leq \ell(\lambda)} \left ( q^{(\lambda_i - \lambda_j) - (i - j)} - 1 \right )}{\prod_{i=1}^{\ell(\lambda)} \prod_{j=1}^{\lambda_i + \ell(\lambda) - i} (q^j-1)}.
\end{equation}

\begin{theorem}[Green {\cite[Theorem~14]{green}}]
\label{degrees}
For all \(n \in \N\), prime powers \(q\), and $\ulambda : \CF (q) \to \Par$ with $\| \ulambda \| = n$,
the degree of the irreducible character $\chi^\ulambda$ of $\GL_n \F_q$ is given by
\begin{equation}
\deg \chi^\ulambda = (q-1)^n \cdot [n]_q! \cdot \prod_{f \in \CF (q)} \left [ \ulambda (f) : q^{\deg f} \right ].
\end{equation}
\end{theorem}

\begin{example}
We use Theorem~\ref{degrees} to calculate $\deg \chi^{f_1 \mapsto (2,1)}$.
By Theorem~\ref{degrees},
\[
\deg \chi^{f_1 \mapsto (2,1)}
=
[3]_2! \cdot [(2,1): 2]
=
[3]_2! \cdot \frac{6}{\prod_{j=1}^3 (2^j-1) \cdot \prod_{k=1}^1 (2^k-1)} = 6,
\]
which agrees with \(\chi^{f \mapsto (2,1)}\) evaluated at \(U_1\) as recorded in Table~\ref{GL3F2_char_table}.
\end{example}



The degrees of primary characters indexed by \(z-1 \mapsto \lambda\) for \(\lambda \vdash n\) have an alternate, combinatorial description, which we mention briefly.
Recall that, for a \myemph{standard Young tableau} \(T\) of shape \(\lambda\), the \myemph{major index} of \(T\) is defined as the sum of all entries \(i\) in \(T\) for which \(i + 1\) appears in a lower row of \(T\) than \(i\).
Furthermore, the major index of \(T\) is denoted \(\maj T\).
Recall also that, given a cell \(a\) in row \(i\) and columns \(j\) of the Ferrers diagram of \(\lambda\), denoted \(a \in \lambda\), the \myemph{hooklength} of \(a\) is defined by \(h(a) = (\lambda_i - i) + (\lambda_j' - j) + 1\).
\begin{theorem}[Stanley \cite{EC2}, Steinberg \cite{steinberg}]
For all \(n \in \N\), \(\lambda \vdash n\), and prime powers~\(q\), we have
\begin{equation}
\deg \chi^{z-1 \mapsto \lambda} = q^{b(\lambda)} \cdot \frac{\prod_{i=1}^n \, q^i - 1}{\prod_{a \in \lambda} \, q^{h(a)} - 1} = \sum_{T} q^{\maj T},
\end{equation}
where the sum ranges over all standard Young tableaux \(T\) of shape \(\lambda\).
\end{theorem}


















\subsection{Certain character values}

Regular semisimple elements have many nice properties. The following theorem, which follows from Steinberg's work, describes one such property.
\begin{theorem}[Steinberg \cite{steinberg}]
\label{sn_stuff}
For all \(n \in \N\), prime powers \(q\), partitions \(\lambda, \mu \vdash n\), and \(g \in \CT_\mu^\Box (q)\), we have \(\chi^{z-1 \mapsto \lambda} (g) = \chi^\lambda_\mu\).
\end{theorem}

Regular elliptic elements, in particular, have nice character-theoretic properties as well.
The following result echoes Corollary~\ref{MNrule_ncycle}.
\begin{corollary}[to~
Corollary~\ref{MNrule_ncycle},
Proposition~\ref{reg_ell_char},
Theorem~\ref{greens_main_thm}, and
Theorem~\ref{degrees}]
\label{reg_ell_chi}
Suppose \(n \in \N\), \(q\) is a prime power, \(\ulambda : \CF(q) \to \Par\) with \(\|\ulambda\| = n\), and \(g \in \CT_{(n)} (q)\).
If \(\chi^\ulambda (g) \neq 0\), then there exist \(d \mid n\), \(f \in \CF_d (q)\), and \(r \in \{0,\ldots,n/d-1\}\) such that \(\ulambda = f \mapsto (n/d - r, 1^r)\) is primary.
Moreover,
\begin{equation}
\deg \chi^{f \mapsto (n/d-r,1^r)} = \deg_{n,d,r} (q),
\end{equation}
as defined in Table~\ref{formula_table}.
\end{corollary}

It follows that, when using the Frobenius formula to enumerate factorizations involving regular elliptic elements, one only needs to consider characters of the form \(\chi^{f \mapsto (n/d-r,1^r)}\).
Therefore, when \(n\) is understood from context, and for any \(d \mid n\), \(f \in \CF_d (q)\), and \(r \in \{0,\ldots,n/d-1\}\), we define
\begin{equation}
\chi^{f,r} = \chi^{f \mapsto (n/d-r,1^r)}.
\end{equation}
We can now write down a further simplified version of~\eqref{ff_in_situ}.
\begin{corollary}[to~Corollary~\ref{ff_in_situ_corollary}~and~Corollary~\ref{reg_ell_chi}]
\label{simpler_ff_in_situ_corollary}
For all \(n,k \in \N\), \(\mu \vdash n\), and prime powers~\(q\), we have
\begin{equation}
\label{simpler_ff_in_situ}
g_{k,\mu} (q)
=
\frac{1}{\gamma_n(q)}
\sum_{d,f,r}
\deg_{n,d,r}(q)^{1-k}
\left ( \sum_{g \in \CT_{(n)} (q)} \chi^{f,r}(g) \right )^k
\left ( \sum_{h \in \CT_\mu (q)}   \chi^{f,r}(h) \right ),
\end{equation}
where the sum is over all \(d\) dividing \(n\), \(f \in \CF_d (q)\), and \(r \in \{0,\ldots,n/d-1\}\).
Moreover, the same is true when both \(g_{k,\mu} (q)\) is replaced with \(g_{k,\mu}^\Box (q)\) and \(\CT_\mu (q)\) is replaced with~\(\CT_\mu^\Box (q)\).
\end{corollary}














































